\section{Conclusion}
We presented OpenMLE, an open full-stack technical solution for training and deploying language-model agents that construct and iteratively improve machine learning solutions through executable feedback. \openmlegym{} provides quality-gated tasks, isolated execution, and task-specific evaluation; \openmleerl{} learns \textsc{Draft}, \textsc{Improve}, \textsc{Debug}, and \textsc{Crossover} transformations through execution-grounded supervised fine-tuning and reinforcement learning; and \openmleevo{} composes the same operators into long-horizon search using structured experience, multi-factor parent selection, and operator-conditioned memory. This shared operator and execution interface makes \thirtyfivebmodel{} both the product of the training stack and the variation engine of its evolutionary harness.

The results show that model learning and search provide complementary gains. Under the identical \openmleevo{} harness, \thirtyfivebmodel{} raises Medal Average over its Qwen3.6-35B-A3B base from $39.39\%$ to \frontisThirtyFiveEvoResult{}, while the companion \thirtybmodel{} reproduces the gain on a second backbone and scale. Matched comparisons further show that \openmleevo{} outperforms general-purpose coding-agent scaffolds across four frontier models and original AIRA-Evo on \thirtyfivebmodel{}; combining \thirtyfivebmodel{} with \openmleevomax{} reaches \frontisThirtyFiveEvoMaxResult{}. Mechanism analyses show sustained late-horizon gains from refinement and recombination, while the matched AIRA-Evo comparison records shorter bounded contexts, lower token use, and more validation progress per token. Controlled comparisons on the ten-task NatureBench Lite subset provide initial evidence that both the post-trained model and the adapted search framework transfer beyond competition-style MLE. Together with the released stack, these results establish OpenMLE as a reproducible testbed for meta-evolution in executable AI4AI.




\section{Authors}

\begingroup
\newcommand{\affmark}[1]{\textsuperscript{#1}}
\setlength{\tabcolsep}{10pt}
\renewcommand{\arraystretch}{1.4}

\begin{tabularx}{\linewidth}{
@{}
>{\raggedright\arraybackslash}X
>{\raggedright\arraybackslash}X
>{\raggedright\arraybackslash}X
@{}
}
Junlin Yang\affmark{1,2,*,$\dagger$}
& Che Jiang\affmark{1,2,*,$\dagger$}
& Yu Fu\affmark{1,3,*} \\

Tianwei Luo\affmark{2,*}
& Can Ren\affmark{1,*}
& Weizhi Wang\affmark{1,2,*} \\

Kaikai Zhao\affmark{2,*}
& Hongyi Liu\affmark{2}
& Yuxin Zuo\affmark{2} \\

Yuru Wang\affmark{1,2}
& Yuchen Fan\affmark{4}
& Kai Tian\affmark{1,2} \\

Zhenzhao Yuan\affmark{1,2}
& Xiaojian Lin\affmark{2}
& Li Sheng\affmark{2} \\

Rushi Qiang\affmark{5}
& Guoli Jia\affmark{2}
& Xingtai Lv\affmark{1,2} \\

Ermo Hua\affmark{2}
& Dianqiao Lei\affmark{1,2}
& Youbang Sun\affmark{2} \\

Ning Ding\affmark{2}
& Bowen Zhou\affmark{2}
& Kaiyan Zhang\affmark{1,$\ddagger$}
\end{tabularx}

\vspace{0.8em}

{\footnotesize\color{frontisgray}
\setlength{\parskip}{2pt}

\textsuperscript{1} Horizon Research, Frontis.AI \quad
\textsuperscript{2} Tsinghua University\quad
\textsuperscript{3} Zhejiang University\par
\textsuperscript{4} Shanghai Jiao Tong University\quad
\textsuperscript{5} Georgia Institute of Technology\par

\vspace{0.35em}

\textsuperscript{*} Core Contributor
\quad
\textsuperscript{$\dagger$} Project Leader\quad
\textsuperscript{$\ddagger$}
Corresponding Author
}
\endgroup

